\section{Results}
\label{sec:results}

\subsection{CNNDetect Results Analysis}
CNNDetect shows a clear domain-shift effect between GAN-era forensic detection and modern diffusion-generated images. In the in-distribution zero-shot setting, the pretrained CNNDetect model achieved only 33.78\% accuracy, 54.60\% ROC AUC, and 25.75\% macro F1-score (Table~\ref{tab:appendix-id-zero-shot}). The confusion matrix shows that the model classified almost all images as real, correctly detecting only 4 fake images out of 600. This resulted in 100.00\% real recall but only 0.67\% fake recall. Although fake precision reached 100.00\%, this value is misleading because the model predicted the fake class only four times.

After fine-tuning on the in-distribution training set, CNNDetect improved substantially. It reached 87.00\% accuracy, 93.80\% ROC AUC, and 84.97\% macro F1-score (Table~\ref{tab:appendix-id-fine-tuned}). The class-level results also became much more balanced, with 79.44\% real F1-score and 90.50\% fake F1-score. This large improvement indicates that the ResNet-based architecture can adapt effectively to Stable Diffusion and Midjourney artifacts when trained on examples from the target distribution.

However, the out-of-distribution results show that this adaptation does not fully transfer to unseen generators. In the OOD zero-shot setting, CNNDetect collapsed almost completely, achieving 33.04\% accuracy, 27.54\% ROC AUC, and 24.83\% macro F1-score (Table~\ref{tab:appendix-ood-zero-shot}). The model failed to correctly classify any fake image, producing 0.00\% fake precision, recall, and F1-score. This confirms that the pretrained detector does not provide reliable zero-shot generalization to fake images generated by unseen diffusion models.

Fine-tuning partially improved OOD performance, but the model still remained weak under generator shift. The fine-tuned CNNDetect achieved 60.52\% accuracy, 57.90\% ROC AUC, and 53.61\% macro F1-score on the OOD set (Table~\ref{tab:appendix-ood-fine-tuned}). It detected fake images better than real images, with 71.51\% fake F1-score compared with 35.71\% real F1-score. This suggests that the fine-tuned model became biased toward predicting the fake class on unseen data.

Overall, CNNDetect demonstrates strong improvement after fine-tuning on known diffusion generators, but its poor zero-shot and reduced OOD performance show that GAN-era forensic features do not generalize reliably to modern diffusion-generated images. The results support the need for detector adaptation and broader generator diversity during training.

\subsection{EfficientNet Results Analysis}

EfficientNet-B4, included as a general-purpose visual classification baseline, showed limited zero-shot capability, but responded well to fine-tuning on the in-distribution data, while failing to generalize to unseen generators.
 
In the \textbf{zero-shot setting}, the ImageNet-pretrained EfficientNet-B4 achieved 50.11\% accuracy, an ROC AUC of 52.21\%, and a macro F1-score of 49.42\%. These results are near random chance, confirming that general-purpose visual representations trained on natural image classification do not transfer to the task of detecting diffusion-generated images without adaptation. 


After \textbf{fine-tuning} on the in-distribution training split, EfficientNet-B4 improved substantially, reaching 88.78\% accuracy, 95.12\% ROC AUC, and 87.49\% macro F1-score. The improvement was consistent across both classes: real F1 rose to 83.47\% and fake F1 to 91.51\%, indicating that the model learned a balanced and reliable decision boundary on the seen generator families. This places EfficientNet-B4 as the second strongest fine-tuned model on the in-distribution test set, behind CLIP+LoRA (95.11\% accuracy, 99.21\% ROC AUC).
 
However, \textbf{out-of-distribution} evaluation revealed that this adaptation did not generalize beyond the training generators. In the zero-shot OOD setting, EfficientNet-B4 scored 50.15\% accuracy, 50.31\% ROC AUC, and 48.09\% macro F1-score — again near chance. After fine-tuning, OOD performance did not recover: accuracy dropped further to 48.22\%, ROC AUC to 48.00\%, and macro F1 to 46.98\%. The fine-tuned model performed worse than random on unseen generators, indicating that the features it learned are specific to Midjourney and Stable Diffusion artifacts present in the training data and do not generalize to the broader space of diffusion generators.


Taken together, these results suggest that EfficientNet-B4 is capable of learning effective in-distribution detection boundaries when fine-tuned, but its convolutional architecture leads it to overfit to the visual fingerprints of the generators seen during training. This makes it unsuitable as a robust detector in open-world settings where new generators are regularly introduced.
 

\subsection{CLIP+LoRA Results Analysis}
CLIP+LoRA achieved the strongest overall performance among all evaluated models, particularly after fine-tuning. Unlike CNN-based forensic detectors that rely heavily on generator-specific artifacts, CLIP begins with a broad visual representation learned from large-scale image-text alignment. The LoRA adaptation mechanism then allows this representation to be specialized to the real-versus-fake classification task while updating only a very small fraction of the total model parameters.

In the zero-shot in-distribution setting, CLIP Zero-Shot achieved 60.67\% accuracy, 50.86\% ROC AUC, and 50.62\% macro F1-score (Table~\ref{tab:appendix-id-zero-shot}). Although the accuracy appears higher than several other zero-shot baselines, the class-level metrics reveal an imbalanced behavior: the model achieved strong fake recall (79.33\%) but weak real recall (23.33\%). This indicates that prompt-based image-text similarity alone does not provide a sufficiently reliable decision boundary for deepfake detection on diffusion-generated images.

After fine-tuning with LoRA adapters, CLIP+LoRA achieved the best in-distribution performance of all evaluated models, reaching 95.11\% accuracy, 99.21\% ROC AUC, and 94.34\% macro F1-score (Table~\ref{tab:appendix-id-fine-tuned}). Performance improved substantially for both classes, with real F1-score increasing to 92.25\% and fake F1-score to 96.43\%. These results indicate that CLIP's pretrained representation can be adapted very effectively to diffusion-generated image detection using lightweight parameter-efficient fine-tuning.

The out-of-distribution results further highlight the strengths and limitations of the CLIP-based approach. In the OOD zero-shot setting, CLIP Zero-Shot achieved 43.38\% accuracy, 52.08\% ROC AUC, and 36.42\% macro F1-score (Table~\ref{tab:appendix-ood-zero-shot}). Similar to the in-distribution zero-shot case, the model strongly favored the fake class, resulting in very low real recall (9.09\%).

After fine-tuning, CLIP+LoRA preserved relatively strong separability under generator shift, achieving 83.45\% ROC AUC on the OOD dataset (Table~\ref{tab:appendix-ood-fine-tuned}), which is substantially higher than the other evaluated models. However, thresholded classification performance remained weaker, with 59.08\% accuracy and 56.12\% macro F1-score. The class-level metrics show that the model correctly identified most fake images (98.11\% fake recall) while misclassifying many real images as fake (29.23\% real recall). This suggests that CLIP+LoRA retains useful ranking capability under distribution shift, but its learned decision threshold does not transfer reliably to unseen generators.

Overall, the CLIP+LoRA results suggest that large pretrained visual representations provide a stronger foundation for diffusion-generated image detection than traditional forensic artifact detectors. Fine-tuning with lightweight LoRA adapters allows the model to adapt effectively while preserving much of CLIP's generalization capability. Among all evaluated models, CLIP+LoRA demonstrated the strongest balance between in-distribution performance and out-of-distribution robustness.


\subsection{DIRE Results Analysis}
The DIRE model shows a clear bias toward detecting fake images across all settings, it predicts nearly every image as fake regardless of content, achieving 66.67\% accuracy that simply mirrors the 2:1 class imbalance, with an AUC of 51.74\% indicating random detection (Table~\ref{tab:appendix-id-zero-shot}).The same behavior appears in the out-of-distribution zero-shot results (Table~\ref{tab:appendix-ood-zero-shot}), which means the model depends on specific patterns from the generators and cannot properly recognize real images. 
After fine-tuning, the model becomes more balanced on in-distribution data (Table~\ref{tab:appendix-id-fine-tuned}),with AUC gaining +18.89 percentage points over the zero-shot, indicating that the model genuinely learned to separate real from fake images, but it is still weaker than other models.  
However, OOD performance collapses to 49.11\% accuracy and an AUC of 33.54\% (Table~\ref{tab:appendix-ood-fine-tuned}). These results contrast strongly with the near-perfect accuracy reported in the original paper \cite{wang2023dire}, which trained and tested on bedroom and ImageNet images from unconditional diffusion models using a matched ADM reconstruction model. Our dataset uses diverse images from modern text-to-image generators like Stable Diffusion and Midjourney, whose architectures produce reconstruction signatures that the ADM model cannot reliably distinguish, breaking the core assumption behind DIRE and explaining both its zero-shot failure and its inability to generalize to unseen generators after fine-tuning. 

\subsection{FIRE Results Analysis}
In the zero-shot in-distribution setting (Table~\ref{tab:appendix-id-zero-shot}), FIRE achieved 66.44\% accuracy, 50.08\% AUC, and 79.73\% F1-score, which indicates weak detection performance despite the high F1-score. The near-random AUC suggests that the model struggled to separate real and fake images and was biased toward predicting the fake class. Compared with the original FIRE paper \cite{chen2025fire}, the performance is lower because the paper used a more controlled setup within the same LSUN Bedroom domain, while this work used the more diverse OpenFake dataset, creating a larger domain gap and making zero-shot transfer more difficult.

After fine-tuning, FIRE improved clearly on the in-distribution dataset (Table~\ref{tab:appendix-id-fine-tuned}), reaching 68.89\% accuracy, 73.20\% AUC, and 77.01\% F1-score. The large increase in AUC from 50.08\% to 73.20\% shows that the model adapted successfully to the OpenFake data. In fact, the achieved AUC is comparable to the range reported in the original FIRE paper \cite{chen2025fire} (around 66–78\%) for challenging generators such as Midjourney and SDXL, even though the current setup is more difficult due to additional dataset variability. This suggests that fine-tuning helped FIRE learn the characteristics of the known generators used in the IOD set.

For OOD evaluation, both zero-shot and fine-tuned FIRE showed poor generalization. Zero-shot OOD (Table~\ref{tab:appendix-ood-zero-shot}) produced 46.66\% AUC, while fine-tuning (Table~\ref{tab:appendix-ood-fine-tuned}) reduced performance further to 43.37\% AUC and 52.67\% accuracy. Although fine-tuning improved in-distribution results, it harmed OOD robustness, indicating that the model became more specialized to the generators seen during training rather than learning generalized fake-image features. Compared with the FIRE paper\cite{chen2025fire}, which focused on a smaller domain shift, the current OOD setting is more challenging because it contains unseen generators and broader distribution differences. Overall, FIRE works reasonably well for known generators after fine-tuning but remains weak in cross-generator generalization.